\documentclass[11pt,a4paper]{amsart}

\usepackage{cgeometry}

\addbibresource{ComplexGeometry.bib}
\defbibheading{bibliography}[\refname]{\section*{#1}}
\setlength{\bibitemsep}{0pt}
\renewcommand*{\bibfont}{\footnotesize}

\allowdisplaybreaks[2]
\emergencystretch=1em

%%%%%%%%%%%%%%%%%%%%%%%%%%%%%%%%%%%%%%%%%%%%%%%%%%%%%%%%%%%%%%%%%%%%%%%%%%%%%
%% Notation
%%%%%%%%%%%%%%%%%%%%%%%%%%%%%%%%%%%%%%%%%%%%%%%%%%%%%%%%%%%%%%%%%%%%%%%%%%%%%

\title[First-stage CLD collapse]{First-stage collapse of the CLD hierarchy on good analytic spaces and arbitrarily branching star products}
\author{}
\date{23 September 2026}

\begin{document}

\begin{abstract}
On a good Berkovich analytic space, every function in the plurisubharmonic sheaf generated by the Chambert--Loir--Ducros decreasing-limit construction is locally a decreasing limit of ordinary smooth plurisubharmonic functions on one fixed ordinary open neighborhood. The same conclusion holds at the tuple of central nodes of every finite product of CLD stars of arbitrary valence. The proofs combine convex averaging of local rows with finite positive barriers. They allow discontinuous limits and changes of tropical presentation between approximants.
\end{abstract}

\maketitle

\section{Introduction}
\label{sec:introduction}

Let $K$ be a complete non-Archimedean field and let $X$ be a topologically separated $K$-analytic space. Chambert--Loir and Ducros give $X$ an affine $G$-tropical structure and define a sheaf $\operatorname{Psh}_X$ by starting with ordinary smooth real functions whose smooth $(1,1)$-form $\mathrm{d}'\mathrm{d}''u$ is positive, then repeatedly taking local decreasing sequential limits \cite[\S12.4.6, Definition~8.1.1 and Lemma~8.1.4]{CLD12}. A section obtained after one such limit is called locally approachable \cite[Definition~8.1.13]{CLD12}. The distinction between a single local approximation and membership in the sheaf generated by transfinite repetition matters when the local smooth rows live on different domains.

The first theorem identifies these two classes when every point has an affinoid analytic domain that is an ordinary neighborhood. Further applications treat finite and arbitrary branching-star products at their node tuples. No single tropical chart is fixed for the approximating sequence in either result.

\begin{theorem}
\label{thm:analytic-first-stage}
Let $X$ be a good, topologically separated Berkovich $K$-analytic space with its Chambert--Loir--Ducros affine structure. For every ordinal $\alpha\geq1$,
\[
\operatorname{Psh}_{X,\alpha}=\operatorname{Psh}_{X,1}=\operatorname{Psh}_X.
\]
Equivalently, every germ of a section of $\operatorname{Psh}_X$ is the pointwise limit of a decreasing sequence of ordinary smooth functions $u_j$ with $\mathrm{d}'\mathrm{d}''u_j\geq0$ on one fixed ordinary open neighborhood of the germ. The limit may take the value $-\infty$.
\end{theorem}

The first theorem concerns the analytic affine structure of \cite[\S12.4.6]{CLD12}. The product theorem uses the affine structure of \cite[\S9.1]{CLD12}. The finite barrier constructions below use features of those structures and do not assert first-stage collapse for every affine $G$-tropical space.

There are two compactness issues in the proof. First, local decreasing rows with a discontinuous common limit need not become uniformly close on overlaps at the same index. A common finite convex average of indices does become uniformly close; this is a weak-convergence argument in a finite product of $C(K)$ spaces. Second, smooth cutoffs are not themselves plurisubharmonic. On an affinoid spectrum, finitely many cutoffs can be controlled by one convex function of truncated logarithms of global analytic functions. The resulting adjusted cutoffs are smooth plurisubharmonic. A regularized maximum then glues the averaged rows without changing their limit. These constructions use the form-level convex composition rule of \cite[Proposition~8.2.1]{CLD12} and the vanishing of $\mathrm{d}'\mathrm{d}''\log|a|$ on an analytic unit locus \cite[Lemma~16.7.4]{CLD12}.

\section{Convex synchronization and diagonalization}
\label{sec:synchronization}

We first record two compactness arguments for sheaves of continuous functions. Neither uses analytic geometry. In the first lemma, all coefficients in the convex average are shared by the different rows.

\begin{lemma}
\label{lem:convex-synchronization}
Let $K_1,\ldots,K_m$ be compact Hausdorff subspaces of a Hausdorff space. For each $i$, let $(f_{i,k})_{k\geq1}$ be a decreasing sequence in $C(K_i,\mathbb{R})$. Suppose these sequences converge pointwise to the restrictions of one function $v$ on $\bigcup_i K_i$, and suppose $v$ is bounded below. For every $\varepsilon>0$ and $k_0\geq1$, there is a finite set of indices $k\geq k_0$ and nonnegative numbers $\lambda_k$ with $\sum_k\lambda_k=1$ such that, with $F_i=\sum_k\lambda_k f_{i,k}$,
\[
\lVert F_i-F_j\rVert_{K_i\cap K_j}<\varepsilon
\qquad\text{for all }i,j.
\]
Moreover $v\leq F_i\leq f_{i,k_0}$ on $K_i$.
\end{lemma}

\begin{proof}
Empty overlaps impose no condition. On the remaining overlaps form the finite product Banach space
\[
B=\prod_{i<j,\ K_i\cap K_j\ne\varnothing} C(K_i\cap K_j,\mathbb{R})
\]
with the maximum norm, and let $b_k$ have $(i,j)$ component $(f_{i,k}-f_{j,k})|_{K_i\cap K_j}$. The sequence $(b_k)$ is uniformly bounded: each $f_{i,k}$ lies between the common lower bound for $v$ and the bounded function $f_{i,1}$. Every component tends pointwise to zero. By the Riesz--Markov representation of the dual of $C(H,\mathbb{R})$ for compact Hausdorff $H$ \cite[Theorem~18.1]{SchenkerFA}, bounded convergence against each finite signed regular measure gives $b_k\rightharpoonup0$ in $B$.

Fix $k_0$. If zero were outside the norm closure of the convex hull of $\{b_k:k\geq k_0\}$, Hahn--Banach separation \cite[Theorem~2.1]{ButtenschoenHB} would give a bounded real linear functional that is uniformly positive on that hull. Its values on the weakly null tail $b_k$ tend to zero, a contradiction. A finite convex combination of tail vectors therefore has norm below $\varepsilon$. Its coefficients give the asserted $F_i$. Monotonicity of each row gives $v\leq F_i\leq f_{i,k_0}$.
\end{proof}

The next lemma replaces a decreasing array of approximants, already defined on one domain, by one decreasing sequence. It uses upper semicontinuity rather than uniform convergence to the intermediate limits.

\begin{lemma}
\label{lem:fixed-domain-diagonal}
Let $K$ be a compact subset of a topological space, and let $\mathscr{S}$ be a class of continuous real functions stable under restriction and addition of constants. For every $j$, suppose $f_{j,k}\in\mathscr{S}(W_j)$ decreases in $k$ to $u_j$ on an open neighborhood $W_j$ of the same $K$. Suppose the functions $u_j|_K$ decrease in $j$ to $u$ on $K$, with values in $\mathbb{R}\cup\{-\infty\}$. Then $u|_{\operatorname{int}K}$ is the decreasing limit of a sequence in $\mathscr{S}(\operatorname{int}K)$.
\end{lemma}

\begin{proof}
On $K$ set $h_n=\min\{f_{j,k}:1\leq j,k\leq n\}$. These functions are continuous, though they need not belong to $\mathscr{S}$, and $h_n\downarrow u$. We use the following compactness fact: if upper semicontinuous $a_r$ decreases to $a$ on $K$ and $b$ is continuous with $a<b$ pointwise, then $a_r<b$ on all of $K$ for some $r$. The open sets $\{a_r<b\}$ increase and cover $K$, so one of them covers $K$.

Set $j_1=k_1=1$ and $F_1=f_{1,1}|_K$. Given $F_{n-1}$ on $K$, put $b_n=\min(h_n,F_{n-1})$. Since $u\leq b_n$, the compactness fact first supplies $j_n$ with $u_{j_n}<b_n+2^{-n-1}$ on $K$, and then supplies $k_n$ with $F_n:=f_{j_n,k_n}|_K<b_n+2^{-n}$ there. Thus $g_n=(f_{j_n,k_n}+2^{-n})|_{\operatorname{int}K}$ belongs to $\mathscr{S}$ and, for $n\geq2$,
\[
g_n\leq F_{n-1}+2^{1-n}=g_{n-1}.
\]
Also $u\leq g_n\leq h_n+2^{1-n}$, so $g_n\downarrow u$ pointwise, including on the set where $u=-\infty$.
\end{proof}

\section{Finite logarithmic barriers}
\label{sec:barriers}

Let $Y=\mathcal{M}(A)$ be a $K$-affinoid spectrum. Write $|a|(y)$ for the seminorm of $a\in A$ at $y$. A logarithmic expression below is evaluated only where the relevant seminorms are positive. When a scalar function is constant for sufficiently negative arguments, its composition with $\log|a|$ has the specified constant extension across $\{a=0\}$.

\begin{lemma}
\label{lem:affinoid-barriers}
Let $K_0\subset Y$ be compact and let $O_1,\ldots,O_m$ be ordinary open subsets covering $K_0$. There are finitely many smooth functions $\theta_1,\ldots,\theta_N:Y\to[0,1]$, each assigned to one of the $O_i$, and a nonnegative smooth function $q:Y\to\mathbb{R}$ such that
\begin{enumerate}
\item the support of each $\theta_j$ lies in its assigned $O_i$, and the ordinary open regions on which some $\theta_j$ is identically $1$ cover $K_0$;
\item $q$, $q+2\theta_j$, and $q-2\theta_j$ are ordinary smooth CLD functions with nonnegative smooth $\mathrm{d}'\mathrm{d}''$ on $Y$ for every $j$.
\end{enumerate}
The functions are built from finitely many elements of $A$, although the number and choice of those elements depend on the cover.
\end{lemma}

\begin{proof}
The evaluation functions $|a|$, $a\in A$, generate the compact Hausdorff topology of $Y$. For each $y\in K_0$ choose a finite system of evaluation inequalities defining an ordinary neighborhood inside one $O_i$, and choose a smaller closed evaluation neighborhood still containing an ordinary neighborhood of $y$. Smooth cutoffs in the finitely many logarithmic variables give a product $\theta$ equal to $1$ near $y$ and supported in the chosen $O_i$. At a coordinate with value zero, use a cutoff equal to $1$ for sufficiently negative logarithmic values and equal to $0$ above a finite threshold. At a positive coordinate, use an ordinary compactly supported cutoff in the logarithmic variable. Each factor is constant at both tails, so the product extends continuously and smoothly across all zero loci. Compactness of $K_0$ yields finitely many products $\theta_i$ whose plateau regions cover it.

List all logarithmic coordinate slots used in these finitely many products. Write $\Theta_i(t_1,\ldots,t_N)$ for the resulting smooth functions on $\mathbb{R}^N$. For each slot $j$ choose a nonnegative smooth compactly supported function $\psi_j$ that is at least $1$ wherever a first or second derivative of the corresponding cutoff factor is nonzero. Put
\[
\rho_j(t)=\int_{-\infty}^{t}(t-s)\psi_j(s)\,\mathrm{d}s.
\]
Then $\rho_j$ is smooth, nonnegative, convex, nondecreasing, and zero on a negative tail. A row and column of $D^2\Theta_i$ vanish whenever the corresponding cutoff factor has zero first and second derivatives. Since the finitely many Hessians $D^2\Theta_i$ are uniformly bounded, one constant $D>0$ can be chosen such that
\[
D\operatorname{diag}(\psi_1,\ldots,\psi_N)\pm 2D^2\Theta_i\geq0
\qquad\text{for every }i\text{ and }t\in\mathbb{R}^N.
\]
Consequently $D\sum_j\rho_j$ and $D\sum_j\rho_j\pm2\Theta_i$ are smooth convex functions on $\mathbb{R}^N$. Compose them with the logarithmic evaluations to define $q$ and $q\pm2\theta_i$.

At a point where some underlying $a_j$ vanishes, its $\rho_j$ summand is identically zero on an ordinary neighborhood and every cutoff factor involving that slot is constant there. Omit those slots. All remaining $a_j$ are analytic units on a smaller ordinary neighborhood: a Laurent domain $|a_j|\geq\varepsilon<|a_j(y)|$ supplies the inverse. Each displayed function is therefore locally a smooth convex function of logarithms of analytic units. The identity $\mathrm{d}'\mathrm{d}''\log|a_j|=0$ as a smooth form on the unit locus \cite[Lemma~16.7.4]{CLD12} and the chain rule show that its smooth Hessian is positive. This includes points on the analytic boundary of $Y$ because the argument concerns smooth forms, not integration by parts.
\end{proof}

\section{Gluing descending rows}
\label{sec:gluing}

For a sheaf $\mathscr{S}$ of continuous real functions, let $\mathscr{A}_0=\mathscr{S}$, let $\mathscr{A}_{\alpha+1}$ be the sheaf of local decreasing limits of $\mathscr{A}_\alpha$ sections, and at a limit ordinal let $\mathscr{A}_\lambda$ be the local union of the earlier stages. Values $-\infty$ are allowed in the latter stages. We prove the gluing statement in the form needed for the CLD generator.

\begin{proposition}
\label{prop:affinoid-gluing}
Let $Y=\mathcal{M}(A)$, and let $\mathscr{S}$ be a sheaf of continuous real functions on $Y$ containing constants and closed under smooth convex coordinatewise nondecreasing composition of finitely many sections. Suppose the functions $q$ and $q\pm2\theta_i$ furnished by \cref{lem:affinoid-barriers} belong to $\mathscr{S}(Y)$ for every finite cover. If $v\in\mathscr{A}_1(U)$ and $K_0\subset U$ is compact, there is an ordinary open neighborhood $W$ of $K_0$ inside $U$ on which $v$ is a decreasing limit of functions in $\mathscr{S}(W)$. Consequently $\mathscr{A}_2=\mathscr{A}_1$ and every later stage agrees with them.
\end{proposition}

\begin{proof}
First suppose $v$ is bounded below on $U$. Choose finitely many open sets $O_i\subset U$ covering $K_0$ and decreasing rows $f_{i,k}\in\mathscr{S}(O_i)$ with limit $v|_{O_i}$. Choose cutoffs and $q$ by \cref{lem:affinoid-barriers}, then reindex the finite cutoff family and repeat its assigned local row for each cutoff. Denote the resulting data again by $O_i,f_{i,k},\theta_i$. By compact Hausdorff regularity, choose compact $K_i$ with
\[
\operatorname{supp}\theta_i\subset\operatorname{int}K_i\subset K_i\subset O_i,
\]
and set $U_i=\operatorname{int}K_i$. Let $V$ be the union of the open plateau regions of the $\theta_i$. Choose a compact neighborhood $L$ of $K_0$ inside $V$. The sets $K_i$, $V$, $L$, and $q$ will remain fixed as the rows are advanced.

Choose an even smooth probability density $\rho_a$ supported in $[-a,a]$. For $r$ real arguments define
\[
R_a^r(t)=\int_{\mathbb{R}^r}\max_{1\leq i\leq r}(t_i+s_i)\prod_{i=1}^r\rho_a(s_i)\,\mathrm{d}s_i.
\]
This function is smooth, convex, nondecreasing in each argument, and translation equivariant. Symmetry of $\rho_a$ and Jensen's inequality give $\max_i t_i\leq R_a^r(t)\leq\max_i t_i+a$. An argument more than $2a$ below another one can be deleted without changing $R_a^r$.

Fix $\eta>0$. Apply \cref{lem:convex-synchronization} to the rows on the compact $K_i$, using a chosen common tail, to obtain finite convex averages $F_i\in\mathscr{S}(O_i)$ such that $|F_i-F_j|<\eta/4$ on $K_i\cap K_j$. Put $p_i=F_i+\eta(2\theta_i-1)$ on $U_i$. At $z\in V$, apply $R_{\eta/16}$ to the arguments $p_i(z)$ for which $z\in U_i$, and call the result $P(z)$. Set $G=P+\eta q$.

We check that $G\in\mathscr{S}(V)$ although the set of available arguments varies. Near each $z\in V$ there is an index $j$ with $\theta_j=1$. If $\theta_i\leq1/4$ at a point of $U_i\cap U_j$, the overlap estimate gives $p_j-p_i>2\eta/16$, so the $i$th argument is inactive. If the boundary of $U_i$ meets $z$, compact support of $\theta_i$ inside $U_i$ makes $\theta_i=0$ near $z$; this argument is inactive wherever it is defined. All other relevant arguments are defined on one neighborhood of $z$. Translation equivariance rewrites $G$ there as the regularized maximum of
\[
F_i+\eta(q+2\theta_i-1).
\]
Each of these is in $\mathscr{S}$ by the hypotheses, and the regularized maximum is an allowed composition. The local formulas agree with the defining formula for $G$, proving membership in the sheaf. The plateau argument and the upper bound for $R_a^r$ also give
\begin{equation}
\label{eq:gluing-bounds}
v+\eta\leq G\leq\max_{i:z\in U_i}F_i+B\eta
\qquad(z\in V),\qquad B=\frac{17}{16}+\sup_Yq.
\end{equation}

Construct $G_n$ by taking $\eta_n=2^{-n}/(B+1)$ and convex averages from row indices at least $n$. For $n\geq2$, choose the tail still farther so that $f_{i,k}<G_{n-1}+\eta_n$ on $K_i\cap L$ for every selected index $k$. Such a common tail exists: the continuous function $G_{n-1}$ strictly exceeds $v$ on $L$ by \eqref{eq:gluing-bounds}, and the decreasing row $f_{i,k}$ converges pointwise to $v$; the increasing open sets $\{f_{i,k}<G_{n-1}+\eta_n\}$ cover the compact $K_i\cap L$. Synchronization can be applied to any prescribed tail. The upper bound in \eqref{eq:gluing-bounds} now yields
\[
G_n\leq G_{n-1}+2^{-n}\quad\text{on }L.
\]
For every $z\in V$, the same bound and $F_i\leq f_{i,n}$ give $G_n(z)\to v(z)$, since only finitely many indices occur. Hence $(G_n+2^{-n})|_{\operatorname{int}L}$ decreases to $v$ and consists of $\mathscr{S}$ sections on one neighborhood of $K_0$.

For arbitrary $v$, put $v^{(m)}=\max(v,-m)$. On $O_i$ the functions
\[
f_{i,k}^{(m)}=\frac{f_{i,k}-m+\sqrt{(f_{i,k}+m)^2+k^{-2}}}{2}
\]
belong to $\mathscr{S}$ by convex monotone composition and decrease to $v^{(m)}$. They are bounded below. Apply the preceding construction to these rows. Its cutoffs, compact sets, and neighborhood $\operatorname{int}L$ depend on the $O_i$ but not on $m$; therefore every $v^{(m)}$ has a decreasing $\mathscr{S}$ row on the same open $\operatorname{int}L$. Since $v^{(m)}\downarrow v$, \cref{lem:fixed-domain-diagonal} on a smaller compact neighborhood of $K_0$ inside $\operatorname{int}L$ gives one decreasing $\mathscr{S}$ sequence converging to $v$ near $K_0$.

For the stage assertion, let $u_j\in\mathscr{A}_1(U)$ decrease to $u$, and fix a compact neighborhood $K_0\subset U$ of a point. Apply the just-proved property to each $u_j$ with this same $K_0$. Each resulting row is defined on a possibly different neighborhood of $K_0$, exactly as allowed in \cref{lem:fixed-domain-diagonal}. The lemma gives a decreasing $\mathscr{S}$ sequence converging to $u$ on $\operatorname{int}K_0$. Thus $\mathscr{A}_2\subseteq\mathscr{A}_1$. The reverse inclusion is immediate, and transfinite induction gives equality at all later stages.
\end{proof}

\section{The analytic CLD sheaf}
\label{sec:cld-sheaf}

\begin{proof}[Proof of \cref{thm:analytic-first-stage}]
On an affinoid analytic space $Y=\mathcal{M}(A)$, let $\mathscr{S}$ be the actual sheaf of ordinary smooth CLD functions with positive smooth $\mathrm{d}'\mathrm{d}''$. Smooth CLD functions form a sheaf of continuous real functions on the ordinary topology \cite[\S6.2.1--6.2.2]{CLD12}. The stage-zero case of the convex composition theorem \cite[Proposition~8.2.1]{CLD12} gives the closure property in \cref{prop:affinoid-gluing}. The functions in \cref{lem:affinoid-barriers} supply its other hypothesis. That proposition shows that every locally approachable function on $Y$ has a decreasing smooth psh approximation on a fixed ordinary neighborhood of any prescribed compact subset, and that the hierarchy stabilizes at stage one. This hierarchy is exactly the one in \cite[Definition~8.1.1 and Lemma~8.1.4]{CLD12}.

For $x\in X$, goodness gives an affinoid analytic domain $Y\subset X$ containing an ordinary open neighborhood $V$ of $x$. The induced analytic structures on $V$ agree, so the ordinary smooth functions, forms, and form positivity on $V$ agree. Their stage-zero sheaves agree, and the sheafified successor and limit operations commute with restriction to $V$. Transfinite induction therefore identifies the corresponding stages on $V$. The affinoid conclusion supplies one ordinary open $W\subset V$ containing $x$ and one decreasing sequence of ordinary smooth psh functions on $W$ converging to the prescribed generated germ. Equality of the stage-one and generated sheaves follows stalkwise.
\end{proof}

\begin{remark}
\label{rem:scope}
The barrier in \cref{lem:affinoid-barriers} uses finitely many global elements of one affinoid algebra to control finitely many changing local presentations. An arbitrary affine $G$-tropical space need not have an affinoid neighborhood or an analogous family of global evaluation functions. The separate criterion in \cref{sec:double-star} uses a common family of smooth affine coordinates.
\end{remark}

\section{A finite double-star product}
\label{sec:double-star}

The gluing argument also applies when ordinary topology is generated on one neighborhood by smooth affine functions. This criterion does not require an analytic realization.

\begin{proposition}
\label{prop:affine-coordinate-criterion}
Let $B$ be a locally compact Hausdorff ordinary open subset of an affine $G$-tropical CLD space. Suppose a family $\mathscr{H}$ of globally defined smooth affine real functions on $B$ generates its ordinary topology. Then the CLD decreasing-limit hierarchy on $B$ stabilizes at stage one. Every generated germ has a decreasing sequence of ordinary smooth psh functions on one fixed ordinary neighborhood.
\end{proposition}

\begin{proof}
Let $K\subset B$ be compact and let finitely many ordinary open row domains cover it. By local compactness, choose for each $x\in K$ a relatively compact ordinary neighborhood whose closure lies in one row domain. Finite-coordinate cylinders in $\mathscr{H}$ form a basis for the ordinary topology, so a smaller such cylinder lies inside this neighborhood. Products of smooth one-variable cutoffs in its finitely many affine coordinates give a smooth function $\theta_x$ equal to $1$ near $x$, with compact support in the assigned row domain. Finitely many of these plateau regions cover $K$.

Collect the finitely many affine coordinates used by these cutoffs. Apply the diagonal Hessian construction in the proof of \cref{lem:affinoid-barriers} to their Euclidean cutoff functions. It gives a nonnegative smooth $q$ such that $q$ and $q\pm2\theta_x$ are smooth convex functions of these affine coordinates. Since each coordinate has zero smooth $\mathrm{d}'\mathrm{d}''$, the convex chain rule makes these functions ordinary smooth psh. The union of the plateau regions lies in the finite union of the compact cutoff supports, so $q$ is bounded on the region used for gluing.

These are precisely the barrier and compactness properties used in \cref{prop:affinoid-gluing}. Its convex synchronization, localizing smooth maximum, truncation, and diagonal argument therefore apply to $B$ with the same estimates. They give a decreasing smooth psh row on one neighborhood of $K$ for every stage-one section and imply equality of all later stages with stage one.
\end{proof}

For a finite set $B_m$ of cardinality $m\geq3$, let $X_{B_m}$ denote the CLD star with branches $[-\infty,0]$ indexed by $B_m$ and with their endpoints at $0$ identified to a node $\eta_m$ \cite[\S6.7.3]{CLD12}. Give products of these stars the induced CLD product affine structure \cite[\S\S6.1.3 and~9.1]{CLD12}.

\begin{theorem}
\label{thm:double-star}
Let $m,n\geq3$. Every germ at $(\eta_m,\eta_n)$ in the generated CLD psh sheaf on $X_{B_m}\times X_{B_n}$ belongs to the first decreasing-limit stage. In particular it has a decreasing sequence of ordinary smooth psh approximants on one fixed ordinary neighborhood of the double node.
\end{theorem}

\begin{proof}
Choose $R,S>0$. In $X_{B_m}$ let $U_R$ be the union of the branch intervals $(-R,0]$, and define $V_S\subset X_{B_n}$ similarly. Because the branch sets are finite, $U_R$ and $V_S$ are ordinary open neighborhoods of their nodes. Enumerate $B_m$ and put
\[
v_1=e_1,\ldots,v_{m-1}=e_{m-1},\qquad v_m=-\sum_{j=1}^{m-1}e_j
\quad\text{in }\mathbb{R}^{m-1}.
\]
On branch $b$, parametrized by $t\in(-R,0]$, define $H_j(b,t)=t(v_b)_j$. Each $H_j$ is smooth affine on $U_R$: it is affine on each branch, and its node slopes sum to zero, as required in the star affine structure \cite[\S6.7.3]{CLD12}. Thus $\mathrm{d}'\mathrm{d}''H_j=0$ as a smooth form by the affine form calculus \cite[\S3.3.11]{CLD12}. The map $H=(H_1,\ldots,H_{m-1})$ embeds $U_R$ into $\mathbb{R}^{m-1}$, since its distinct branch rays meet only at the origin. Indeed, it is the restriction of the continuous injection from the compact truncated star with intervals $[-R,0]$, hence a topological embedding. Construct an analogous embedding $G$ on $V_S$.

The functions $H_j\circ\mathrm{pr}_1$ and $G_\ell\circ\mathrm{pr}_2$ are smooth affine on $U_R\times V_S$ by the CLD product affine rule \cite[\S9.1.4]{CLD12}. They have zero smooth Hessian and together generate the ordinary product topology because $(H,G)$ is a topological embedding. Apply \cref{prop:affine-coordinate-criterion} on this ordinary open neighborhood of $(\eta_m,\eta_n)$.
\end{proof}

The finite-coordinate construction in \cref{thm:double-star} uses a neighborhood from which the $-\infty$ endpoints have been removed. It gives no assertion for arbitrary nonambient branching systems whose ordinary neighborhoods lack a common family of topology-generating affine functions.

\section{Products of arbitrarily branching stars}
\label{sec:arbitrary-star-products}

\begin{proposition}
\label{prop:positive-coordinate-criterion}
Let $B$ be a locally compact Hausdorff ordinary open subset of an affine $G$-tropical space. Suppose its ordinary topology is generated by a family of globally defined ordinary smooth psh functions. Then the decreasing-limit hierarchy on $B$ stabilizes at stage one.
\end{proposition}

\begin{proof}
Fix a compact $K\subset B$ and a finite ordinary open cover on whose members a stage-one section has decreasing smooth psh rows. Choose finitely many smooth cutoffs $\theta_i$ that equal $1$ on ordinary open sets covering $K$ and have compact support in their assigned row domains. They can be taken as smooth functions $\Theta_i(h_1,\ldots,h_N)$ of finitely many topology-generating smooth psh functions: take a finite-coordinate cylinder around each point with closure in an assigned row domain, apply compactly supported Euclidean cutoff functions, and collect the finitely many coordinates used. After collecting them, multiply each cutoff by extra factors in the other coordinates, chosen equal to $1$ near its plateau point. This keeps its support in the assigned domain and makes each $\Theta_i$ compactly supported on $\mathbb{R}^N$.

The first and second derivatives of all $\Theta_i$ are supported in a common compact subset of $\mathbb{R}^N$. Choose $D>0$ such that the functions
\[
Q(t)=D\sum_{j=1}^N\mathrm{e}^{t_j},\qquad Q(t)\pm2\Theta_i(t)
\]
are convex and nondecreasing in each coordinate on $\mathbb{R}^N$. On the compact derivative support, this follows by making $D\mathrm{e}^{t_j}$ dominate the finitely many first derivatives and the operator norms of the Hessians; outside it the derivatives of $\Theta_i$ vanish. Hence $q=Q(h_1,\ldots,h_N)$ and $q\pm2\theta_i$ are smooth psh by the stage-zero convex monotone composition rule \cite[Proposition~8.2.1]{CLD12}. The function $q$ is bounded on the compact region used for gluing. The row synchronization and gluing proof of \cref{prop:affinoid-gluing} now applies on that region, including the truncation and diagonal arguments. Every stage-two germ is stage one, and transfinite induction gives the assertion.
\end{proof}

For an arbitrary set $B$ with at least two elements, let $X_B$ be the CLD star of \cite[\S6.7.3]{CLD12} and let $o$ be its node. Fix a branch $b_*$ and $R>0$, and remove the closed tail $\{b_*\}\times[-\infty,-R]$. The complement $U$ is an ordinary open node neighborhood: all other branches remain complete. For each $b\neq b_*$ and $a<0$, choose a smooth convex nondecreasing function $\rho_a$ that vanishes on $(-\infty,a]$, is strictly increasing on $(a,0]$, and is affine of slope $1$ near $0$. Put $c_a=\rho_a(0)$ and define $h_{b,a}$ on $U$ by
\[
h_{b,a}(b,t)=\rho_a(t),\qquad h_{b,a}(b_*,t)=c_a-t,
\qquad h_{b,a}=c_a\quad\text{on all other branches}.
\]
It is ordinary smooth: it is constant near every included $-\infty$ endpoint, changes on only two branches, and its node slopes $1$ and $-1$ sum to zero. Its smooth Hessian vanishes near the node and on the reservoir branch; on the transition part of branch $b$ it equals $\rho_a''(t)\,\mathrm{d}'t\wedge\mathrm{d}''t\geq0$. Thus $h_{b,a}$ is stage zero.

These functions generate the ordinary topology of $U$. On a nonreservoir branch, $h_{b,a}<c_a$ isolates that branch and varying $a$ separates its points and endpoint neighborhoods. On the reservoir branch, $h_{b,a}>c_a$ isolates it and gives a strictly monotone coordinate. At the node, a basic ordinary neighborhood excludes tails in only finitely many branches. For each such nonreservoir branch one inequality in $h_{b,a}$ removes its tail; one further inequality removes any prescribed reservoir tail. Their finite intersection lies inside the basic neighborhood. Continuity gives the reverse inclusion of the generated topology.

\begin{theorem}
\label{thm:arbitrary-star-products}
Let $X_{B_1},\ldots,X_{B_r}$ be CLD stars with arbitrary branch sets $B_i$ of cardinality at least two, where $r$ is finite. At the tuple of central nodes in their product, every germ of the generated CLD psh sheaf belongs to the first decreasing-limit stage.
\end{theorem}

\begin{proof}
In each factor choose the ordinary open neighborhood $U_i$ and smooth psh family $\{h_{b,a}\}$ above. Pullbacks of these functions by the product projections remain smooth psh by \cite[Proposition~8.1.5 and \S9.1.4]{CLD12}. Their union generates the ordinary product topology on $\prod_i U_i$. Apply \cref{prop:positive-coordinate-criterion} there and restrict to the node tuple.
\end{proof}

\printbibliography

\end{document}
